\section*{Chapter \ref{ch:qm:linear-time-series} --- Linear Time Series}

\begin{solution}{exo:qm:linear-time-series:1}
\omterm{def:qm:stochastic-differential-equations:ou}{Half-life} $\ln(1/2)/\ln0.98 = 34.3$ days; $\rho(20) = 0.98^{20} = 0.67$.
\end{solution}

\begin{solution}{exo:qm:linear-time-series:2}
No: the root of $1 + 2z$ is $-1/2$, inside the unit circle. The MA(1) with coefficient $1/2$ and innovation variance four times larger, $X_t = \eta_t +
\frac12\eta_{t-1}$, has the same autocovariances ($\rho(1) = \vartheta/(1 + \vartheta^2) = 0.4$ in both cases) and is invertible.
\end{solution}

\begin{solution}{exo:qm:linear-time-series:3}
$w_0 = 1$, $w_1 = -d = -0.3$, $w_2 = w_1(1 - d)/2 = -0.105$, $w_3 = w_2(2 - d)/3 = -0.0595$.
\end{solution}

\begin{solution}{exo:qm:linear-time-series:4}
$-(1 + 3 \times 0.99)/1\,000 = -0.0040$: the estimate averages 0.986 when the truth is 0.99. The \omterm{def:qm:stochastic-differential-equations:ou}{half-life} at 0.99 is 69.0 days; a naive analyst who
estimates 0.99 should correct to 0.994, a \omterm{def:qm:stochastic-differential-equations:ou}{half-life} of 114.6 days. A bias of 0.4\% in $\phi$ is a bias of 66\% in the \omterm{def:qm:stochastic-differential-equations:ou}{half-life}.
\end{solution}

\begin{solution}{exo:qm:linear-time-series:5}
Multiply $X_t = \phi_1X_{t-1} + \phi_2X_{t-2} + \varepsilon_t$ (demeaned) by $X_{t-h}$, $h \ge 1$, and take expectations: $\gamma(h) = \phi_1\gamma(h - 1) + \phi_2\gamma(h - 2)$;
divide by $\gamma(0)$. At $h = 1$, $\rho(1) = \phi_1 + \phi_2\rho(1)$, so $\rho(1) = 0.5/1.3 = 0.385$; then $\rho(2) = 0.5 \times 0.385 - 0.3 = -0.108$.
\end{solution}

\begin{solution}{exo:qm:linear-time-series:6}
$T = 100$: $-2.86154 - 0.028903 - 0.000423 - 0.000040 = -2.891$. $T = 12\,574$: $-2.862$.
\end{solution}

\begin{solution}{exo:qm:linear-time-series:7}
\texttt{spurious()}: $|t| > 1.96$ in 89\% of the 5\,000 regressions of levels, with median $R^2$ 0.17; regressing the differences, in 5.0\%, the nominal
rate. Differencing removes the common stochastic trend problem when the series are not cointegrated (chapter~20).
\end{solution}

\begin{solution}{exo:qm:linear-time-series:8}
The $t$-statistic of the lagged level in that regression is the Dickey--Fuller $\tau$, whose 1\% and 5\% critical values are $-3.43$ and $-2.86$: $-2.79$
is not significant at 5\%. The \omterm{def:qm:stochastic-differential-equations:ou}{half-life} of 552 days is also biased down (739 after Kendall's correction) and very uncertain (a 10--90\% range of about
350 to 970 days if the truth were 739). A trade that waits years for reversion must be sized for the possibility that there is none.
\end{solution}

\begin{solution}{pb:qm:linear-time-series:1}
\textbf{1.} Mean 84.5 bp; range $-241$ to 291 bp; 25 bp on 22 September 2026.
\textbf{2.} 0.9987; the correlogram declines almost linearly, still 0.64 after 250 days, the signature of a \omterm{def:qm:linear-time-series:unitroot}{unit root} or near \omterm{def:qm:linear-time-series:unitroot}{unit root}.
\textbf{3.} 4.6 bp; 0.015, 0.015, $-0.027$.
\textbf{4.} AR(10), with every coefficient below 0.035 in absolute value: detectable in 12\,574 days, worthless for prediction.
\textbf{5.} Approximately \omterm{def:qm:linear-time-series:stationary}{white noise}: $\Delta X_t \approx \varepsilon_t$ with $\psi_j$ close to zero for $j \ge 1$.
\textbf{6.} $\hat\phi = 0.99875$, \omterm{def:qm:estimation:estimator}{standard error} 0.00045, \omterm{def:qm:stochastic-differential-equations:ou}{half-life} 552 trading days (2.2 years).
\textbf{7.} $t = -2.79$, and a normal one-sided $p$-value of 0.003.
\textbf{8.} $(1 + 3\hat\phi)/n = 0.00032$, so $\phi = 0.99906$ and a \omterm{def:qm:stochastic-differential-equations:ou}{half-life} of 739 days (2.9 years).
\textbf{9.} With 0, 5, 10 and 20 lags: $-2.79$, $-2.74$, $-3.13$, $-3.55$, against critical values $-3.43$ (1\%) and $-2.86$ (5\%). Not rejected at 5\% with
few lags, rejected at 5\% with ten and at 1\% with twenty; AIC picks 34 lags and $-3.14$.
\textbf{10.} $\tau = -2.01$: no evidence against a \omterm{def:qm:linear-time-series:unitroot}{unit root}.
\textbf{11.} Median 576 days, 10--90\% range 352 to 969.
\textbf{12.} In 57\% of the histories.
\textbf{13.} $H = 0.34$ for the changes, $d = -0.16 \pm 0.06$: over-differenced, so the level has $d \approx 0.84$, a slowly and hyperbolically mean-reverting
process.
\textbf{14.} Because the true memory is not geometric, the fitted AR(1) coefficient depends on the sample's length and period; and near one the \omterm{def:qm:stochastic-differential-equations:ou}{half-life}
magnifies every error.
\textbf{15.} The distribution of the path to reversion: how far the spread can move against the position and for how long, which depends on the volatility
and on the tails of the memory, not on the \omterm{def:qm:stochastic-differential-equations:ou}{half-life} alone.
\textbf{16.} The data cannot say: consistent with a very persistent \omterm{def:qm:linear-time-series:stationary}{stationary process} and with a \omterm{def:qm:linear-time-series:unitroot}{unit root}.
\textbf{17.} The bias-corrected 739 days, treated as uncertain between one and four years, with the position sized to survive no reversion at all.
\textbf{18.} Under a \omterm{def:qm:linear-time-series:unitroot}{unit root} the $t$-statistic follows the Dickey--Fuller law, whose 5\% point is $-2.86$, not $-1.645$; the true $p$-value is about 0.06.
\textbf{19.} Named result: \emph{the \omterm{def:qm:stochastic-differential-equations:ou}{half-life} of a spread}: the AR(1) \omterm{def:qm:stochastic-differential-equations:ou}{half-life} of the 2s10s spread is 552 trading days, 739 after Kendall's
bias correction; the Dickey--Fuller statistic $-2.79$ does not reject a \omterm{def:qm:linear-time-series:unitroot}{unit root} at 5\%, and augmented versions reject or not depending on the lag choice.
\textbf{20.} Whether the data rule out a random walk, which, for a spread one hopes will revert, is usually the answer that it has not been ruled out.
\end{solution}

\begin{solution}{iq:qm:linear-time-series:1}
$\ln(1/2)/\ln\phi \approx 0.69/(1 - \phi)$ periods. Near one, a small error in $\hat\phi$ is a large relative error in $1 - \hat\phi$, the estimate is biased down by about
$(1 + 3\phi)/n$, and its distribution is skewed; the \omterm{def:qm:stochastic-differential-equations:ou}{half-life} can be off by a factor of two.
\iqlookfor{The formula, the $1/(1 - \phi)$ sensitivity, and the small-sample bias.}
\end{solution}

\begin{solution}{iq:qm:linear-time-series:2}
Under $\phi = 1$ the regressor is a random walk; the \omterm{def:qm:estimation:estimator}{estimator} converges at rate $n$, not $\sqrt n$, and its $t$-statistic converges to a functional of \omterm{def:qm:brownian-motion:bm}{Brownian
motion} (the Dickey--Fuller law), shifted to the left: 5\% critical value $-2.86$ with a constant, instead of $-1.645$.
\iqlookfor{Nonstandard asymptotics and the critical values.}
\end{solution}

\begin{solution}{iq:qm:linear-time-series:3}
Whether it is a \omterm{def:qm:linear-time-series:unitroot}{spurious regression}: are both series unit-root processes, and are the residuals stationary (a cointegration test, chapter~20)? Regress
changes on changes, or test the residuals with Engle--Granger critical values; look at the residuals' autocorrelation.
\iqlookfor{\omterm{def:qm:linear-time-series:unitroot}{Spurious regression} as the first hypothesis, and cointegration as the test.}
\end{solution}

\begin{solution}{iq:qm:linear-time-series:4}
An AR($p$) has a \omterm{def:qm:linear-time-series:acf}{partial autocorrelation function} that cuts off after lag $p$ and an \omterm{def:qm:linear-time-series:acf}{autocorrelation function} that decays; an MA($q$) has the reverse, an
\omterm{def:qm:linear-time-series:acf}{autocorrelation function} that cuts off after lag $q$. ARMA has both decaying; use an \omterm{def:qm:linear-time-series:ic}{information criterion} to choose.
\iqlookfor{The cut-off rules and a model-selection tool.}
\end{solution}

\begin{solution}{iq:qm:linear-time-series:5}
Autocorrelations that decay like a power of the lag and do not sum: the fractional $d$ between 0 and one half. In markets: volatility (absolute and squared
returns), trading volume, order-flow signs, some spreads and interest rates; rarely in returns themselves.
\iqlookfor{Power-law decay, and volatility as the canonical example.}
\end{solution}

\begin{solution}{iq:qm:linear-time-series:6}
The regression uses the sample mean, which absorbs part of each deviation, and the \omterm{def:qm:estimation:estimator}{estimator} is a ratio of correlated quadratic forms; both push it down. Kendall's
approximation with an estimated mean is $\E[\hat\phi] - \phi \approx -(1 + 3\phi)/n$: about $-4/n$ near one, so $-0.016$ with a year of daily data.
\iqlookfor{The mechanism and the $-(1 + 3\phi)/n$ size.}
\end{solution}
